% ECONOMICS_PROBLEM: {"id": "C65-45", "title": "Viscosity existence without the BSDE comparison structure", "jel_code": "C65", "source_id": "OP-0563"}
% FIRST_POSTED: 2026-10-09
% ATTEMPT_STATUS: solved
% ENTRY_KIND: research_attempt
% !TeX program = pdflatex
% Self-contained: no external input or bibliography files required.
\documentclass[11pt,a4paper]{article}
\usepackage[T1]{fontenc}
\usepackage[utf8]{inputenc}
\usepackage{lmodern}
\usepackage[margin=24mm]{geometry}
\usepackage{amsmath,amssymb,amsthm,mathtools}
\usepackage{booktabs,longtable,array,enumitem,microtype,xurl,hyperref}
\hypersetup{colorlinks=true,linkcolor=blue,citecolor=blue,urlcolor=blue,pdftitle={Checked research attempts: nine economics and stochastic-analysis problems}}
\setlength{\parindent}{0pt}
\setlength{\parskip}{3pt}
\setlength{\emergencystretch}{3em}
\widowpenalty=10000\clubpenalty=10000\displaywidowpenalty=10000
\setlist[enumerate]{leftmargin=*,itemsep=2pt,topsep=3pt}
\newcommand{\R}{\mathbb R}
\newcommand{\N}{\mathbb N}
\newcommand{\E}{\mathbb E}
\newcommand{\Pp}{\mathbb P}
\newcommand{\Q}{\mathbb Q}
\newcommand{\F}{\mathcal F}
\newcommand{\ind}{\mathbf 1}
\newcommand{\cE}{\mathcal E}
\newcommand{\Law}{\operatorname{Law}}
\newcommand{\sgn}{\operatorname{sgn}}
\newcommand{\Var}{\operatorname{Var}}
\newcommand{\dist}{\operatorname{dist}}
\newcommand{\KL}{\operatorname{KL}}
\newcommand{\TV}{\operatorname{TV}}
\newcommand{\Lip}{\operatorname{Lip}}
\DeclareMathOperator*{\esssup}{ess\,sup}
\DeclareMathOperator*{\argmax}{arg\,max}
\newtheorem{theorem}{Theorem}
\newtheorem{lemma}[theorem]{Lemma}
\newtheorem{proposition}[theorem]{Proposition}
\newtheorem{corollary}[theorem]{Corollary}
\theoremstyle{definition}
\newtheorem{example}[theorem]{Example}
\newcommand{\report}[3]{\clearpage\setcounter{theorem}{0}\renewcommand{\thetheorem}{#1.\arabic{theorem}}\renewcommand{\theHtheorem}{#1.\arabic{theorem}}\section*{#1: #2}\addcontentsline{toc}{section}{#1: #2}\textbf{Outcome:} #3.\par\textbf{Tools:} web browsing available; Python computation available. Literature checked on 9 October 2026.}
\newcommand{\stage}[2]{\subsection*{#1) #2}}
\newcommand{\unresolved}{\textbf{LABEL: UNRESOLVED.}}
\newcommand{\disproof}{\textbf{LABEL: FULL SOLUTION. SUBLABEL: COUNTEREXAMPLE/DISPROOF.}}

\begin{document}
\section*{C65-45: Viscosity existence without comparison structure}
\textbf{Outcome:} FULL SOLUTION: counterexample to the statement as written.\par\textbf{Tools:} web browsing available; Python computation available. Literature checked on 9 October 2026.
\subsection*{1) FORMAL RESTATEMENT}
For every $T>0$, integers $n,d\ge1$, and continuous coefficients $b,\sigma,g,h$ satisfying the two growth inequalities and the spatial Lipschitz assumptions in the uploaded statement, is there a function $u\in C([0,T]\times\R^n)$ with terminal value $h$ that solves
\[
 u_t+b\cdot Du+\tfrac12\operatorname{tr}(\sigma\sigma^\top D^2u)
       +g(t,x,u,\sigma^\top Du)=0
\]
in the stated viscosity sense? At a local maximum of $u-\varphi$ in $(0,T)\times\R^n$, the displayed expression, evaluated on $\varphi$ and $u$, must be nonnegative; at a local minimum it must be nonpositive. Tests are $C^{1,2}$. No growth restriction on the candidate continuous solution is added.

The constants satisfy $K>0$, $k\ge0$, $\alpha\in(1,2]$, $\alpha^*=\alpha/(\alpha-1)$, and $1\le p<\alpha^*$. In particular, $\sigma\equiv0$ is permitted. This is the decisive stress point. The other stress point is a continuous, non-Lipschitz dependence on $y$. We use $\sgn(0)=0$ and do not replace the question by a uniformly elliptic version.

\subsection*{2) QUICK LITERATURE/CONTEXT CHECK}
The primary source is Problem 5.2 of Fan--Hu--Tang \cite{FHT}, with growth conditions (4.13)--(4.14). The PDF numbers its accompanying existence result Theorem 4.10; the linked HTML uses different theorem numbering. The PDF assumptions explicitly permit zero diffusion. The argument below is independent of that existence theorem and makes no claim of bibliographic priority.

\subsection*{3) ATTACK PLAN}
\textbf{Proof routes tested.} A BSDE-based construction would need continuous selection of solutions as the initial state varies. Perron's method would need adequate continuous barriers or a comparison substitute. Neither property follows just from pointwise ODE existence.

\textbf{Disproof routes.} First set $b=\sigma=0$; then force opposite nonunique ODE branches by a sign-changing spatial parameter. Finally check that viscosity solutions cannot evade the ODE obstruction. This is the selected route.

\textbf{Landscape and tools.} This is a degenerate semilinear PDE and selection problem. The relevant tools are time reversal (to fix signs), viscosity penalization (to restrict to a spatial slice), a one-dimensional monotonicity test (to identify viscosity ODE solutions), an integral identity (to avoid an unproved uniqueness theorem), scalar differential inequalities (to force branch separation), and direct growth estimates (to verify every assumption).

\subsection*{4) WORK}
\begin{theorem}[Explicit nonexistence example]
Set
\[
 T=1,\quad n=d=1,\quad b=0,\quad\sigma=0,\quad h=0,
 \qquad g(t,x,y,z)=x+2\sgn(y)\sqrt{|y|}.
\]
These data satisfy every assumption in the question, with
$K=1$, $k=2$, $\alpha=2$, $\alpha^*=2$, and $p=1$.
Nevertheless, there is no continuous viscosity solution on $[0,1]\times\R$.
\end{theorem}
\begin{proof}[Verification of the coefficients]
The map $y\mapsto\sgn(y)\sqrt{|y|}$ is continuous, including at zero. All remaining coefficients are continuous. Both spatial Lipschitz inequalities for $b,\sigma$ hold with left side zero. Write $r=|y|$. For $y\ne0$,
\[
 \sgn(y)g\le |x|+2\sqrt r\le |x|+1+r
       \le2(1+|x|+r+|z|^2).
\]
At $y=0$ the one-sided left side is zero. Further,
\[
 |g|+|h|\le |x|+1+r
       \le2(1+|x|+e^{2r}+|z|^2).
\]
The first inequality uses $(\sqrt r-1)^2\ge0$, and the second follows, for example, from $e^{2r}\ge1+2r$. Thus both required growth inequalities hold globally. Also $p=1<2=\alpha^*$.
\end{proof}

We give the viscosity argument in detail; absence of a classical solution alone would not suffice.

\begin{lemma}[Restriction to spatial slices]
Let $F:\R\times\R\to\R$ be continuous. Suppose $v$ is a continuous viscosity solution of
$v_s-F(x,v)=0$ on $(0,1)\times\R$, with the conventional subsolution inequality $\varphi_s\le F(x,v)$ at a maximum. Then $w(s)=v(s,x_0)$ is a viscosity solution of $w'=F(x_0,w)$ for every $x_0$.
\end{lemma}
\begin{proof}
Let $w-\phi$ have a local maximum at $s_0\in(0,1)$. Replace $\phi$ by $\phi_\eta(s)=\phi(s)+\eta(s-s_0)^2$, $\eta>0$, to make this maximum strict on a sufficiently small closed time interval. Choose a compact spatial interval about $x_0$ and maximize
\[
 v(s,x)-\phi_\eta(s)-N|x-x_0|^2
\]
on the resulting rectangle. Its maximum is at least its value at $(s_0,x_0)$. Boundedness on the rectangle implies $|x_N-x_0|\to0$. Every subsequential limit of $s_N$ maximizes $w-\phi_\eta$ on the time interval; strictness gives $s_N\to s_0$. The maximizing points are therefore interior for all sufficiently large $N$. Adding a constant to the penalized test makes it touch $v$ there. Since the PDE has no spatial derivatives, its subsolution inequality gives
$\phi_\eta'(s_N)\le F(x_N,v(s_N,x_N))$. Continuity and passage to the limit give $\phi'(s_0)\le F(x_0,w(s_0))$. For a local minimum, minimize $v-\phi_\eta+N|x-x_0|^2$, strictifying in the opposite direction. This proves the supersolution inequality as well.
\end{proof}

\begin{lemma}[Continuous viscosity ODE solutions are classical]
If a continuous $w:[0,1)\to\R$ solves $w'=F(x_0,w)$ in viscosity sense on $(0,1)$, then
\[
 w(s)=w(0)+\int_0^s F(x_0,w(r))\,dr.
\]
\end{lemma}
\begin{proof}
Put $A(s)=\int_0^s F(x_0,w(r))dr$ and $r(s)=w(s)-A(s)$. The function $A$ is $C^1$. Adding $A$ to any test for $r$ shows that $r$ solves $r'=0$ in viscosity sense.

Here is a direct proof that a continuous viscosity subsolution of $r'\le0$ is nonincreasing. Suppose $r(b)>r(a)$ for $0<a<b<1$. Choose $\eta>0$ with $\eta(b-a)<r(b)-r(a)$, and choose $c\in(b,1)$. For sufficiently small $\delta>0$,
\[
 r(b)-\eta b-\frac\delta{c-b}
 >r(a)-\eta a-\frac\delta{c-a}.
\]
The function $r(t)-\eta t-\delta/(c-t)$ tends to $-\infty$ as $t\uparrow c$ and has a maximum on $[a,c)$ at some point in $(a,c)$. The corresponding touching test has derivative $\eta+\delta/(c-t)^2>0$, a contradiction. Thus $r$ is nonincreasing. Applying the same argument to $-r$, which is a subsolution when $r$ is a supersolution, makes $r$ nondecreasing. It is constant on $(0,1)$ and, by continuity, at zero too. This proves the integral identity without assuming local Lipschitz continuity of $F$.
\end{proof}

\begin{proof}[Nonexistence]
Assume a continuous viscosity solution $u$ exists and put $v(s,x)=u(1-s,x)$. Time reversal turns precisely the sign convention in the question into
\[
 v_s=F(x,v),\qquad F(x,y)=x+2\sgn(y)\sqrt{|y|},\qquad v(0,x)=0.
\]
Indeed a touching test $\phi(s,x)$ corresponds to $\phi(1-t,x)$ and its time derivative changes sign. The two lemmas show that each slice $w(s)=v(s,x)$ is a classical solution of this ODE, including its right derivative at zero.

Fix $x>0$. Since $w'(0)=x$, $w(s)>0$ for small positive $s$. It cannot first return to zero: on a positive interval $w'=x+2\sqrt w\ge x$, so integration from any positive point makes $w$ strictly increase. Consequently $w(s)>0$ for every $s\in(0,1)$. On this interval $q(s)=\sqrt{w(s)}$ satisfies
\[
 q'(s)=1+\frac{x}{2\sqrt{w(s)}}\ge1.
\]
Integrating from $\varepsilon$ to $s$ and letting $\varepsilon\downarrow0$ gives $q(s)\ge s$, hence $v(s,x)\ge s^2$ for every $x>0$. If $x<0$, then $-w$ solves the positive-parameter equation with parameter $-x$. Therefore $v(s,x)\le-s^2$ for every $x<0$.

For fixed $s\in(0,1)$, continuity at $x=0$ would imply simultaneously
$v(s,0)\ge s^2$ and $v(s,0)\le-s^2$. Since $s^2>0$, this is impossible.
\end{proof}

\subsection*{5) VERIFICATION}
\textbf{Tiny cases.} At $x=0$, the ODE has the solutions $0$, $s^2$, $-s^2$, and delayed variants; selecting any one cannot repair the limits from both spatial sides. At $T=0$ the contradiction disappears, as it must. The example needs only one space coordinate, one Brownian coordinate, and zero terminal data.

\textbf{Computational check.} For $x>0$, separation of variables yields
\[
 s=q-\frac{x}{2}\log(1+2q/x),\qquad q=\sqrt{v(s,x)}.
\]
The supplied script solves this scalar equation by bracketing. At $s=1/2$ and $x=10^{-1},10^{-2},10^{-3},10^{-6}$ it gives respectively
$0.3975850$, $0.2738441$, $0.2534698$, and $0.2500069$ for $v$; the negative side is its negative. These finite tests illustrate, but are not used to prove, the gap of at least $1/2$ between the two sides.

\textbf{Adversarial checks.} No uniqueness of the non-Lipschitz ODE was invoked. Spatial restriction was proved for viscosity tests rather than assumed. Only compact rectangles were used in penalization, so no bounds at spatial infinity are hidden. The absence of diffusion is authorized by the literal assumptions; imposing uniform ellipticity would be a substantive new problem, not a correction of a typo. The contradiction rules out every continuous solution, not merely polynomial-growth solutions or solutions obtained by a particular BSDE selection.

\subsection*{6) FINAL}
\disproof

The explicit data in the theorem satisfy all the uploaded hypotheses and admit no continuous viscosity solution. The complete proof consists of coefficient verification, the two viscosity lemmas, and the forced opposite inequalities $v(s,x)\ge s^2$ for $x>0$ and $v(s,x)\le-s^2$ for $x<0$. No claim is made here about a version that adds uniform ellipticity or a suitable one-sided uniqueness condition in $y$.

\clearpage
\begin{thebibliography}{99}
\bibitem{FHT}
Shengjun Fan, Ying Hu and Shanjian Tang. \emph{1D nonlinear backward stochastic differential equations: a unified theory and applications}. arXiv:2309.06233v2, 11 February 2026. In particular Problems 5.1--5.5, PDF pp.\ 32--33; for C65-45 see (A1), (4.13)--(4.14), PDF p.\ 31. \url{https://arxiv.org/pdf/2309.06233v2}.

\end{thebibliography}
\end{document}
